% New source after checkpoint08. Root reserved A.4.49.
\subsection{Finite fusion of the translation copies with the multiplier retained}

Write $\mathcal R(d)$ for the bound
$I(dD_{289}(d)^{45})^I$ in (ECR.5), where (ECR.2)--(ECR.3)
are started with $D_0=d$.  Retain the actual quotient
$q:\mathcal M\to Q=G^\dagger$ of (ECR.1), put $D=|Q|$, and
apply A.4.47 once to obtain
$\eta_0:\mathcal M\to E_0$ with its projection to $Q$.
Apply its construction a second time, with $\eta_0$ as the
specified old quotient, and include the period coordinate
(ECR.21).  Denote the resulting maps by
\[
 \eta_1^*:\mathcal M\longrightarrow E_1^*
          \longrightarrow E_0\longrightarrow Q.
                                                               \tag{CFF.1}
\]
Each arrow in (CFF.1) is induced by the identity on the original
mapping class group.  Uniform bounds for these particular maps are
\[
 b_0=\mathcal R(D),\qquad b_1=\mathcal R(b_0),\qquad
 |E_0|\leq b_0,\quad |E_1^*|\leq b_*:=24b_1^{25}.
                                                               \tag{CFF.2}
\]
Here the second characteristic-cover recursion may be started
with the bound $b_0$ instead of computing the exact order of $E_0$.
The larger initial value preserves every kernel containment in its
proof.

Use the original surface input and the genus-one compiler bounds
\[
 \begin{gathered}
 s_0=128\cdot3^{1172},\qquad
 m=\max\{b_*,m_{s_0,9}\},\\
 U=\mathcal M/\kappa_m(P\mathcal M),\qquad
 Z=Z_1(m)=\Gamma_{0,5}/\kappa_{N(m)}(\Gamma_{0,5}),
 \end{gathered}                                                \tag{CFF.3}
\]
where $m_{s_0,9}$ and $N(m)$ are the explicit expressions
(FGA.1)--(FGA.3).  Thus the input (CFF.3) includes the finite
same-parameter relation compiled in A.4.39; no full pentagon premise
is being substituted for that finite test.
The maps to $E_1^*$ and $E_0$ factor through $U$ by the actual
regular-permutation argument of A.4.40.  Write these factorizations
as $\pi_1$ and $\pi_0$.

Let $\iota$ be the actual pure orientation-reversing involution
of the surface polygon in (TPL.8).  Its based surface map is
\[
 \iota(A)=B,\qquad\iota(B)=A,\qquad
 \iota(C_i)=(C_1\cdots C_{i-1})C_i^{-1}
                              (C_1\cdots C_{i-1})^{-1}.
                                                               \tag{CFF.4}
\]
The literal free-word verification retained there proves its
square is identity and specifies every peripheral conjugator.
Conjugation by this extended mapping class preserves
$P\mathcal M$ and its characteristic subgroup $\kappa_m(P\mathcal M)$.
It therefore gives an actual automorphism $\jmath$ of $U$.
This use of $\jmath$ occurs before projection in (CFF.1).

\begin{result}{A.4.49}{Proved}{One explicit finite five-point input gives geometric fusion at the recognition level for every multiplier}
For every actual gentle pair $F=(\lambda,f)$ whose pentagon value is
one in the specified group $Z$ of (CFF.3), let $\Psi_F\in\operatorname{Aut}(U)$
be the finite action of A.4.40.  Define $e\in\{0,1\}$ and the corrected
finite source action by
\[
 (-1)^e=\lambda\bmod3,\qquad
 \Psi_F^+=\jmath^e\circ\Psi_F .
                                                               \tag{CFF.5}
\]
There is an actual pure orientation-preserving mapping class
$g\in P\mathcal M$ for which
\[
       \pi_0\bigl(\Psi_F^+(\bar h)\bigr)
                    =\eta_0(g h g^{-1})\qquad(h\in H).
                                                               \tag{CFF.6}
\]
Here $\bar h$ is the original named element in $U$, so (CFF.6)
retains both translation generators elementwise.
Consequently the actual homomorphism
\[
 \beta_F:U\longrightarrow E_0,\qquad
 u\longmapsto\eta_0(g)^{-1}\pi_0(\Psi_F^+(u))\eta_0(g)
                                                               \tag{CFF.7}
\]
fixes the named $H$ and has
\[
 \begin{split}
 p_0\beta_F(\overline{N_{\mathcal M}(H)})
                       &\subset q(N_{\mathcal M}(H)),\\
 p_0\beta_F(\overline{C_{\mathcal M}(H)})
                       &\subset q(C_{\mathcal M}(H)).
 \end{split}                                                   \tag{CFF.8}
\]
The two finite outputs $d\operatorname{pr}_G$ and aligned ordinary
forgetting agree on the zero-label part of the first image in
(CFF.8).  The source of $\beta_F$ is the displayed universal finite
group $U$; no assertion that it descends to an automorphism of
$E_0$ is required.
\end{result}

\begin{proof}
\par\medskip\noindent\textbf{The actual finite complement domain.}\quad
The chosen value of $m$ allows all maps in (CFF.1), the period map
$\Omega$ from (ECR.21), and the surface extraction of A.4.40.
That theorem supplies simultaneously $\Psi_F$ and the actual
peripheral automorphism
\[
 F_\Pi\in\operatorname{Aut}
       \bigl(\Pi/\kappa_{s_0}(\Pi)\bigr),
 \qquad \Pi=\pi_1(S_{1,9}),
                                                               \tag{CFF.9}
\]
with the universal-curve conjugation square on the actual
geometric image.  Its peripheral power is $\lambda$.
The same square for $\Psi_F^+$ uses
$F_\Pi^+=\iota^e F_\Pi$, with peripheral power
$\mu=(-1)^e\lambda$.  In particular $\mu\equiv1\pmod3$.
Both the original action and the actual reflection preserve each
puncture label, so $\Psi_F^+$ is identity on the label permutation
quotient.  The finite forgetting square to the unpunctured torus
is supplied by A.4.40 and the ordinary reflection, with its actual
inner path comparison.  Every $h\in H$ has identity compact period;
its image under that square still has identity period in $\Omega$.
An inner path comparison also sends identity to identity.

It follows that
$h\mapsto\pi_1(\Psi_F^+(\bar h))$ has its image in the exact subgroup
(ECR.22) for this second construction:
\[
 B_1=(K_0/V_1)\rtimes H,\qquad
 V_1=\ker(\eta_1^*|_{K_0}).
                                                               \tag{CFF.10}
\]
Its projection to $H$ is the identity, because the retained labels
are faithful.  Thus it is an actual complement in this specified
finite semidirect product.  The projection of this complement to
$(K_0/V_0)\rtimes H$, where $V_0=\ker(\eta_0|_{K_0})$, is exactly
$h\mapsto\pi_0(\Psi_F^+(\bar h))$.
This proves the source identification needed for the finite-cover
argument; a tuple of matrices with the same orders is not being
used in place of this homomorphism.

\par\medskip\noindent\textbf{A reduced-cone fixed-point argument without a chosen lift.}\quad
Use the $289$ characteristic refinements of (ECR.12) starting at
$V_0$ and ending at $V_1$.  Each finite cover $Y_i$ is connected
and has the actual finite cell model of (ECR.13).  The complement
in (CFF.10), and its reductions, act on these covers: write its
elements as $(k_h,h)$, acting by a deck transformation followed
by the specified lift of the configuration action.  These actions
and the covering maps are equivariant for the same abstract $H$.
They need not fix the selected basepoint.

Put $k=\mathbb F_3$ and use the functorial cochain complex
\[
 \widetilde C^*(Y_i)
       =\operatorname{Cone}\bigl(k\longrightarrow C^*(Y_i;k)\bigr).
                                                               \tag{CFF.11}
\]
The copy of $k$ is in degree minus one, with differential the
inclusion of constant zero-cochains.  Since $Y_i$ is nonempty and
connected, (CFF.11) is equivariantly quasi-isomorphic to the
quotient complex $C^*(Y_i;k)/k$, in degrees zero through sixteen.
Its cohomology is zero in nonpositive degrees and is ordinary
cohomology in positive degrees.  Hence every covering map in the
construction induces zero on all its cohomology groups.

For $A\cong C_3$ in $H$, apply the complete Tate double-complex
argument (ECR.14) to this quotient complex.  Seventeen successive
maps give zero Tate hypercohomology.  The restriction
\[
 \operatorname{Cone}\bigl(k\to C^*(Y_i;k)\bigr)
       \longrightarrow
 \operatorname{Cone}\bigl(k\to C^*(Y_i^A;k)\bigr)
                                                               \tag{CFF.12}
\]
is a Tate isomorphism: the identity on $k$ cancels from its cone,
and the remaining relative cellular complex has only free $A$
orbits.  The bounded-cochain horizontal calculation proving this
is exactly the one in A.4.47.  The action on the right fixed
subcomplex, including its copy of $k$, is trivial for $A$.
Thus zero Tate maps give zero in every ordinary cohomology degree
of the reduced cone on the $A$-fixed subcomplex.  This includes
degree minus one when that fixed subcomplex is empty.

For completeness, the dimension bound for this next cyclic stage
is four.  Choose $a$ generating $A$ and a complementary $b$ in $H$.
A configuration fixed by $A$ satisfies
\[
 z_{h+a}=z_h+t,\qquad 3t=0,\quad t\ne0.
                                                               \tag{CFF.13}
\]
For each of the eight possible values of $t$, the only continuous
parameters are $z_b$ and $z_{2b}$ in the torus, subject to the
distinct-point conditions.  The $A$-fixed locus is therefore a
union of open submanifolds of real dimension four.  In the
corresponding manifold cover $X_i$ of $X$, the fixed set maps
locally homeomorphically onto this fixed locus: on an invariant
evenly covered neighborhood, uniqueness of the lift fixing its
center proves that assertion.  The finite cover $Y_i$ is the
lifted retract subcomplex inside $X_i$; its $A$-fixed cells are
contained in $X_i^A$ and therefore have dimension at most four.
The reduced cone
for this fixed set lies in degrees minus one through four.

Let $Z_j=Y_{17j}^A$.  By (CFF.12), each of the seventeen successive
maps between the $Z_j$ has zero map on all reduced-cone
cohomology.  Apply the complete Tate calculation to the residual
$H/A\cong C_3$ action.  There are only six possible vertical
degrees.  Six maps already give zero on the Tate abutment and
hence on every reduced-cone cohomology group of the $H$-fixed
sets.  The remaining maps preserve that vanishing.  We obtain
\[
 H^{-1}\operatorname{Cone}\bigl(k\to C^*(Y_0^H;k)\bigr)
       \longrightarrow
 H^{-1}\operatorname{Cone}\bigl(k\to C^*(Y_{289}^H;k)\bigr)
                         \quad\hbox{equal to zero}.             \tag{CFF.14}
\]
The bottom fixed set $Y_0^H$ cannot be empty.  Otherwise every
fixed set above it would be empty by equivariance of the covering
maps, and both groups in (CFF.14) would be $k$, with the identity
map induced by constants.  This contradicts (CFF.14).
Notice that this argument proves nonemptiness at the required
lower cover, without presupposing a fixed point at the top.

\par\medskip\noindent\textbf{The forty-eight ordinary sections and their exact markings.}\quad
Choose a fixed point in $Y_0$.  Its image in $X$ is one of the
forty-eight grid embeddings $\xi:H\to T[3]$ of (NCF.4).
Choose once and for all a configuration path $\gamma_\xi$ from
the reference grid to this labelled grid and an isotopy extension
$g_{\gamma_\xi}$ on the torus fixing the zero point.
At its endpoint the underlying marked set is the same $T[3]$.
The endpoint has period identity and label permutation equal to
the matrix $\bar A$ of $\xi$.
Choose an integral matrix $A\in\operatorname{GL}_2(\mathbb Z)$
reducing to $\bar A$ and with determinant
$\epsilon\in\{1,-1\}$ equal to $\det\bar A$ in $\mathbb F_3$.
Such a lift is obtained from the elementary generators of
$\operatorname{SL}_2(\mathbb F_3)$ and the integral reflection
for the other determinant; these generate all forty-eight
matrices.  Let $\ell_A$ be its actual linear torus map.  Then
\[
 \alpha_\xi=g_{\gamma_\xi}^{-1}\ell_A
                                                               \tag{CFF.15}
\]
is a pure extended mapping class, with orientation sign
$\epsilon$ and with each peripheral class sent to its
$\epsilon$ power up to conjugacy.

Here is the corresponding section in the configuration extension.
For a path $\gamma$ from the reference grid to a fixed grid,
the lift fixing its endpoint has, in the path model of the
fundamental extension, the cocycle
\[
        k_h=[\,\gamma*(h\gamma)^{-1}\,],\qquad h\in H.
                                                               \tag{CFF.16}
\]
Indeed $k_h(h\gamma)$ is the path $\gamma$, and
$k_h\,h(k_{h'})=k_{hh'}$ by cancellation of the two middle paths.
The normalized evaluation fibration of (NCF.1) identifies this
section with
$\alpha_\xi h\alpha_\xi^{-1}$: the linear map sends the
translation $T_{-h/3}$ to $T_{-\xi(h)}$, and the isotopy extension
returns its grid along $\gamma_\xi$.  The two evaluation paths
are precisely the two paths in (CFF.16).  Thus this comparison
also retains the negative-translation parameter fixed explicitly
in A.4.46.

The chosen fixed point of $Y_0$ may be a different lift of $\xi$
from the one selected by $\gamma_\xi$.  Select its actual deck
displacement $\bar k\in K_0/V_0$ and an ordinary lift $k\in K_0$.
Uniqueness of the lift fixing a given point gives equality of
the two finite sections after conjugation by $\bar k$.  Therefore
the actual pure extended mapping class
\[
                         \alpha=k\alpha_\xi
                                                               \tag{CFF.17}
\]
satisfies (CFF.6) with $\alpha$ in place of $g$.
This construction gives a finite procedure: enumerate the fixed
vertices in the actual finite cover, choose a path in its
one-skeleton, and use its literal point-pushing word for the deck
displacement.  It retains the displacement instead of setting
an unknown marking path to zero.

\par\medskip\noindent\textbf{The negative grid sign is excluded in the retained surface quotient.}\quad
It remains to prove that $\epsilon=1$ for the reflected action
chosen in (CFF.5).  The universal-curve square (CFF.9) and the
projection of (CFF.17) to the old $G$ factor say exactly that
\[
             F_0=\alpha^{-1}F_\Pi^+
                  \quad\hbox{centralizes the named }H
                    \hbox{ in }\operatorname{Out}(Q_L).
                                                               \tag{CFF.18}
\]
Here the expression uses the based action of the actual extended
mapping class $\alpha$.  All changes of based lift are inner
and hence disappear on the abelianization used next.
Both factors preserve each named peripheral conjugacy class up
to its common power, so $F_0$ has common power
$\mu\epsilon$.  All its required quotients are actual maps from
the characteristic surface input in (CFF.9): the verbal $Q_L$
and the separate pure-peripheral compact quotient $H_c$ are
the maps (JP.1), retained by $s_0$.
There is no requirement that the deck subgroup itself preserve
$H_c$.

In the literal basis $(A,B,C_1,\ldots,C_8)$ of (TFC.3), write
$\delta=\mu\epsilon\bmod3$.  The homology matrix of $F_0$ and
the integral deck matrices reduced modulo three are
\[
 F_0^{\mathrm{ab}}=
       \begin{pmatrix}C&0\\ B&\delta I_8\end{pmatrix},
 \qquad
 h_i^{\mathrm{ab}}=
       \begin{pmatrix}I_2&0\\ t_i&P_i\end{pmatrix}.
                                                               \tag{CFF.19}
\]
Equation (CFF.18) implies the exact thirty-two equations
\[
              (I-P_i)B=t_i(C-\delta I_2),\qquad i=a,b.
                                                               \tag{CFF.20}
\]
This calculation is valid for the actual peripheral transport;
it does not assert that $F_0$ belongs to the geometric image.
Stack the row-major entries of the $a$ equation before those
of the $b$ equation.  On a stacked right side $r\in\mathbb F_3^{32}$
the four functionals
\[
 r_1+r_3,\qquad r_2+r_4,\qquad
                 2r_{21}+2r_{27},\qquad2r_{22}+2r_{28}
                                                               \tag{CFF.21}
\]
annihilate the left side for every $8$-by-$2$ matrix $B$.
Substitution of the two literal matrices (TFC.4) shows that on
the right side they give, in order, the four entries of
$C-\delta I_2$.  Hence (CFF.20) forces $C=\delta I_2$
without any prior determinant assumption.

For the determinant comparison use the separate map to $H_c$.
It kills the eight specified puncture generators, and the
ninth peripheral becomes $[A,B]^{-1}$.  This commutator is
central of order nine in that actual class-two quotient.
The pure peripheral map $F_0$ therefore gives
$[A,B]^{\det C}=[A,B]^{\mu\epsilon}$, and in particular
\[
                     \det C=\delta\pmod3.
                                                               \tag{CFF.22}
\]
But $C=\delta I_2$ gives $\det C=\delta^2=1$.
Thus $\delta=1$.  Since $\mu\equiv1\pmod3$, this proves
$\epsilon=1$.  Consequently (CFF.17) is an actual pure
orientation-preserving mapping class; take $g=\alpha$.
This proves (CFF.6) for every multiplier using exactly the
reflection in (CFF.5).

The executed certificate
\path{work/tower_code/configuration_grid_orientation.py}
rebuilds both integral deck matrices from the actual free
words, verifies the rank-twelve left-side space and rank-sixteen
joint space, and checks all four functionals (CFF.21).
It also checks all eighty-one compact matrices for both signs.
The independent replay retains the same matrices and the
same four functional rows.  These computations concern this
precisely specified orientation obstruction, not a census of
gentle source images.

\par\medskip\noindent\textbf{Projection through the canonical recognition map.}\quad
Equation (CFF.6) proves that the actual homomorphism (CFF.7)
fixes the original named translations.  The image under it of
every ordinary normalizer element therefore normalizes
$\eta_0 H$, and the corresponding assertion holds for
centralizers.  Applying the first construction A.4.47 gives
exactly the inclusions (CFF.8).  On zero-label ordinary
normalizers, (MFP.4) identifies descent with aligned forgetting;
that equality therefore holds on the stated finite images.
Every use of finite invertibility in this construction is the
actual record-power argument of A.4.40.  Neither the proof nor
the conclusion requires the finite pentagon locus to be a group.
\end{proof}

\par\medskip\noindent\textbf{What the explicit level supplies.}\quad
Equations (CFF.1)--(CFF.3) specify a single finite five-point
source and actual quotient maps for this translation fusion and
normalizer comparison.  The result applies to every genuine
gentle source element meeting that finite pentagon test, with
the reflection residue and the pure conjugator retained.
It does not report an executed enumeration of that source's
gentle shadows or of their parity values.
